
\documentclass[a4paper,12pt]{article}

\input{misc/basic.tex}
\usepackage[%
  style=ext-authoryear-tecomp,
  backend=biber,
%  introcite=plain,
  natbib=true,
  citetracker=true,
  maxcitenames=1,
  maxbibnames=1000,
  isbn = false,
  eprint =false,
  doi = true,
  url = false, 
  giveninits=true, 
  citestyle=authoryear,
  uniquename=init,
]{biblatex}

\AtEveryCitekey{\ifciteseen{}{\defcounter{maxnames}{3}}}

\DeclareFieldFormat[inbook,article]{citetitle}{#1} %Entfernt Quotation Marks in Titel
\DeclareFieldFormat[inbook,article]{title}{#1} %Entfernt Quotation Marks in Titel


%%%% ABSTAND LITERATURVERZEICHNIS
\setlength{\bibhang}{10mm}
\setlength{\bibitemsep}{5pt}


%\setlength{\introcitewidth}{10pt}
%\setlength{\introcitesep}{12pt}

\renewbibmacro{in:}{}
%\renewbibmacro{URL:}{}
\renewcommand*{\nameyeardelim}{\addcomma\space}


\DeclareDelimFormat[parencite]{finalnamedelim}{\addspace\&\space}
%\DeclareNameAlias{default}{family-given}
\DeclareNameAlias{sortname}{family-given}



\DefineBibliographyStrings{ngerman}{
   andothers = {{et\,al\adddot}},
%   pages = {},
}

\usepackage{csquotes} % Für Error Message beheben
\input{misc/colors.tex}

%\usepackage{biblatex}
\addbibresource{build/wiar.bib}
\nocite{*}


\usetikzlibrary{arrows.meta}

\usepackage{times}
\usepackage[font={small}]{caption}
\usepackage{lipsum}
\usepackage{fancyhdr}
\usepackage{titletoc}
\usepackage[
    activate={true,nocompatibility},
    final,
    tracking=true,
    kerning=true,
    spacing=true,
    factor=1100,
    stretch=30,
    shrink=20
]{microtype}

\newgeometry{left=2.5cm,right=2.5cm,top=2.5cm,bottom=2.5cm,
		includeheadfoot,
		headheight=15pt
		}
\setlength{\parindent}{1cm}

\renewcommand{\baselinestretch}{1.5}
\newcommand{\absatz}{\par \vspace*{6pt}} %6pt ODER Leerzeile

\titlecontents{section}[2.8em]
		{\mdseries}
		{\contentslabel{2.8em}}
		{}
		{\titlerule*[2mm]{.} \contentspage}
\titlecontents{subsection}[4.2em]
		{\mdseries}
		{\contentslabel{2.8em}}
		{}
		{\titlerule*[2mm]{.} \contentspage}		


\titleformat{\section}[hang]
		{\Large \sffamily \bfseries}
		{\thesection}
		{4mm}
		{}
\titleformat{\subsection}[hang]
		{\large \sffamily \bfseries}
		{\thesubsection}
		{4mm}
		{}

\pagestyle{fancy}
\fancyhf{}
\fancyhead[OR]{\thepage}
\fancyhead[OL]{\leftmark}


\begin{document}
\thispagestyle{empty}



\begin{center}
%\raggedbottom

\includegraphics[height=2cm]{misc/unibonnlogo}
\includegraphics[height=2cm]{misc/ieslogo}
\par \vspace*{2mm}
RHEINISCHE FRIEDRICH-WILHELMS-UNIVERSITÄT BONN 
\par \vspace*{-2mm}
{\small Agrar-, Ernährungs- und Ingenieurwissenschaftliche Fakultät}
\par \vspace*{1.5cm}
{\large \bfseries \sffamily Wissenschaftliche Arbeit}
\par \vspace*{6mm}
im Rahmen des Moduls
\par
 \vspace*{-1mm}
{\bfseries \sffamily
Wissenschaftliches Arbeiten in der Agrar- und Ernährungsökonomie
}
\par \vspace*{2mm}


\par \vspace*{1cm}
{\Large ,,Wie lässt sich die Lagerstabilität \par \vspace*{-4mm} von Mikronährstoffen \par \vspace*{-4mm} in angereicherten Backwaren \par unter tropischen Bedingungen absichern?''}
\par 
 – Eine systematische Literaturübersicht –
 \par \vspace*{1.5cm}
{\footnotesize vorgelegt von:}
\par \vspace*{0mm}
Leander Ihle - 50156139 \par \vspace*{-2mm}
Samuel Knorr - 50147935 \par \vspace*{-2mm}
Lars Wombacher - 50167241
\par \vspace*{10mm}
{\footnotesize vorgelegt am:}
\par \vspace*{0mm}
\today
\par 
\vfill
1. Prüfende: Chiara Backes
\par \vspace*{-2mm}
2. Prüfende: Prof. Dr. Denise Fischer-Kreer



\end{center}





\pagebreak
\thispagestyle{empty}
\newgeometry{left=2.5cm,right=2.5cm,top=2.5cm,bottom=2.5cm,
%		includeheadfoot,
		headheight=15pt
		}
\tableofcontents


\pagebreak
\newgeometry{left=3cm,right=2.5cm,top=2.5cm,bottom=2.5cm,
		includeheadfoot,
		headheight=15pt
		}
\setcounter{page}{1}


\section{Einführung}
Das Bewusstsein für den Erhalt an Mikronährstoffen in Lebensmitteln über eine längere Lagerungszeit spielt eine immer wichtiger werdende Rolle in der Lebensmittelforschung. Diese reichen von Lebensmittelversorgung in Kriegsgebieten, Gebieten von Naturkatastrophen bis hin zur Lebensmittelversorgung in Luft- und Raumfahrt. \par 

Diese wissenschaftliche Arbeit bearbeitet die Fragestellung nach einem möglichen Mikronährstoffabbau in angereicherten Backwaren unter tropischen Bedingungen. Ziel ist es, mehr Klarheit für Produzenten zu schaffen, die eine Gewährleistung von Mindestmengen an Nährstoffen in ihren Produkten benötigen.  \par 

So hat man bereits festgestellt, dass Nudeln angereichert mit Spinat- und Karottenpulver auch noch nach 4 Monaten grundsätzlich ausreichende Mikronährstoffgehalte vorweisen konnten \autocite{sharma}. Daten dieser Studie zeigten aber auch, dass unter tropischen Bedingungen ein allgemeiner Mikronährstoffverlust von bis zu 40\% zu verzeichnen gewesen war. Nicht zu vernachlässigen, gilt hier aber auch das Packungsmaterial, welches neben der Temperatur einen weiteren großen Effekt ausmacht. 


\par \vspace*{4mm}

\pagebreak
\section{Methoden}
\subsection{Suchstrategie}
Um die Forschungsfrage zu beantworten wurde eine systematische Literaturanalyse durchgeführt und dabei die PRISMA Richtlinien dazu befolgt. Die Literaturrecherche erfolgte über Google Scholar. Der folgende Suchbegriff wurde angewandt:
\textit{''micronutrient§'' AND ''ambient storage'' AND ''fortified'' -fruits}.
Die Suchergebnisse wurden auf Arbeiten aus den letzten 5 Jahren beschränkt (2021 - 2026). Desweiteren wurden nur Peer-reviewed Studien für die Untersuchung der Forschungsfrage herangezogen. Die einzelnen Schritte der Recherche sind im PRISMA-Flussdiagramm übersichtlich dargestellt.


\subsection{Ein- und Ausschlusskriterien}
Studien wurden einbezogen, wenn diese relevant zur Forschungsfrage waren. Es wurden nur Publikationen mit einbezogen, die auf Deutsch oder Englisch verfasst wurden. Um die große Menge an Studien zu dem Thema einzugrenzen, wurden Studien vernachlässigt, die Lebensmittel untersuchten, welche für die Forschungsfrage zu unterschiedlich waren (wie beispielsweise Früchte). Desweiteren wurden alle Studien, die Themenfremd waren exkludiert.

\subsection{Daten-Management}
Da diese Arbeit von mehreren Autoren bearbeitet wurde, wurde eine Citavi-Cloud angelegt und Textelemente über Sciebo ausgetauscht.

\pagebreak

\begin{figure}[h]
\centering
\resizebox{.8\textwidth}{!}{
\begin{tikzpicture}[>=Latex, font={\sf \small}]

\tikzstyle{bluerect} = [rectangle, rounded corners, minimum width=1.5cm, minimum height=0.75cm, text centered, draw=black, fill=cyan!60!gray!45!white, rotate=90, font=\sffamily]

\tikzstyle{roundedrect} = [rectangle, rounded corners, minimum width=12cm, minimum height=1cm, text centered, draw=black, font=\sffamily]

\tikzstyle{textrect} = [rectangle, minimum width=5.25cm, text width=5.24cm, minimum height=1cm, draw=black, font={\sffamily \footnotesize},inner sep=3mm]

\node (top1) at (0, 10.5cm) [draw, roundedrect, fill=yellow!80!red!70]
  {Identification of studies via database and registers};

\node (r1blue) at (-6.75cm, 8.0cm) [draw, bluerect, minimum width=3.5cm]{Identification};

\node (r1left) at (-3.25cm, 8.0cm) [draw, textrect, minimum height=3.5cm]
  {Volltexte aus Datenbanksuche 'Google Scholar'
		($n=167$)
 
  };

\node (r1right) at (3.25, 8.0cm) [draw, textrect, minimum height=3cm]
  {Volltexte, die vor Screening entfernt wurden
    \begin{itemize}
    \item Duplikate ($n=0$)
    \item Journal-Ranking ($n=114$)
    \item Nicht in deutsch/englisch ($n=2$)
    \end{itemize} 
  };

\node (r2blue) at (-6.75cm, 1.5cm) [draw, bluerect, minimum width=8.5cm]
  {Screening};

\node (r2left) at (-3.25cm, 5.0cm) [draw, textrect, minimum height=1.5cm]
  {Volltexte Titel Screening ($n=51$)};

\node (r2right) at (3.25, 5.0cm) [draw, textrect, minimum height=1.5cm]
  {Ausgeschlossene Volltexte ($n=41$)};

%\node (r3left) at (-3.25cm, 2.5cm) [draw, textrect, minimum height=1.5cm]
%  {Reports sought for retrieval ($n=$)};

%\node (r3right) at (3.25, 2.5cm) [draw, textrect, minimum height=1.5cm]
%  {Not retrieved ($n=$)};


\node (r4left) at (-3.25cm, 0cm) [draw, textrect, minimum height=1.5cm]
  {Reports assessed for eligibility ($n=10$)};

\node (r4right) at (3.25, -1.0cm) [draw, textrect, minimum height=3.5cm]
  {Reports Excluded: 
    \begin{itemize}
    \item Ungeeignetes Lebensmittel ($n=2$)
    \item Forschungsfrage wird nicht bearbeitet ($n=1$)
%    \item Reason 3 ($n=$)
%    \item etc.
    \end{itemize}     
  };


\node (r5blue) at (-6.75cm, -4.5cm) [draw, bluerect, minimum width=2cm]
  {Included};

\node (r5left) at (-3.25cm, -4.5cm) [draw, textrect, minimum height=2cm]
  {Einbezogene Studien \\ 
   $n= 7$ \\ 
%   Reports of included studies \\ 
%   $n=$  
   };


% Draw arrows between nodes:
\draw[thick, ->] (r1left) -- (r1right);
\draw[thick, ->] (r1left) -- (r2left);
\draw[thick, ->] (r2left) -- (r2right);
%\draw[thick, ->] (r2left) -- (r3left);
%\draw[thick, ->] (r3left) -- (r3right);
%\draw[thick, ->] (r3left) -- (r4left);
\draw[thick, ->] (r2left) -- (r4left);
\draw[thick, ->] (r4left) -+ (0.35,0);
\draw[thick, ->] (r4left) -- (r5left);

\end{tikzpicture}
}
\caption{PRISMA-Flussdiagramm}
\end{figure}


\pagebreak
%\nocite{*}
\printbibliography

\pagebreak
\thispagestyle{empty}
\section*{Erklärung}

%Hiermit erkläre ich, dass ich die vorliegende Arbeit selbstständig und ohne Benutzung anderer als der angegebenen Hilfsmittel angefertigt habe. Alle Stellen, die wörtlich oder sinngemäß aus veröffentlichten und nicht veröffentlichten Schriften entnommen
%wurden, sind als solche kenntlich gemacht. Die Arbeit ist in gleicher oder ähnlicher Form oder auszugsweise im Rahmen einer anderen Prüfung noch nicht vorgelegt worden. \par

Hiermit versichern wir (Namen der Studierenden), dass wir die vorliegende Arbeit selbstständig und ohne Inanspruchnahme anderer als der
ausdrücklich angegebenen Hilfsmittel angefertigt haben. Alle Stellen, die wörtlich oder
sinngemäß aus veröffentlichten oder unveröffentlichten Schriften übernommen wurden, sind
als solche gekennzeichnet.
In der vorliegenden Arbeit haben wir KI-basierte Werkzeuge (wie z.B. ChatGPT Version X;
XX) wie folgt verwendet: 
\par \vspace*{4mm}
\begin{multicols}{2}

\begin{itemize}
\item[$\square$] gar nicht
\item zur Ideenfindung
\item zur Strukturierung
\item zur Formulierung einzelner Textpassagen im Umfang von \underline{\hspace*{5mm}}\% des gesamten Textes
\item zur Datenanalyse, falls ja, bitte spezifizieren:
\item zur Entwicklung von Software-Quellcode
\item zum Korrekturlesen oder zur Verbesserung des Textes
\item zum Übersetzen des Textes
\item sonstiges, nämlich:\underline{\hspace*{10mm}}
\end{itemize}
\end{multicols}
\par \vspace*{4mm}
Die Autorinnen/ Autoren übernehmen die volle Verantwortung für den Inhalt der Arbeit. Die
Arbeit wurde nicht in gleicher oder ähnlicher Form oder in Auszügen für eine andere
Prüfungsleistung eingereicht.

\vspace*{4mm}
Bonn, den \today
\par \vspace*{10mm}
 
Unterschrift der/des Studierenden 


\end{document}